% GENERATED FILE. Edit canonical lessons, then run npm run generate:books.
% canonical-source-hash: fnv1a64:b90b95eb0ad1839b
% canonical-lessons: JA-C134-denwa, JA-W134-de, JA-C134-densha, JA-C134-kaban, JA-W134-ba, JA-C134-basho, JA-C134-taberu, JA-W134-be, JA-C134-benkyou, JA-C134-shinbun, JA-W134-bu, JA-R134-phone-and-train

\chapter{Phone and Train: De, Ba, Be and Bu}
\label{ch:ja-phone-and-train-de-ba-be-bu}
\hlchaptermodality{134}

\begin{chapteropening}
\textbf{By the end of this chapter:} \emph{I can write de, ba, be and bu by adding the voicing mark to signs I already write, and read words that use them, such as denwa, densha and shinbun.}

Read denwa, densha, kaban, basho, taberu, benkyou and shinbun, write the four new signs beside the signs they voice, write densha and basho from memory and say which one you can ride, and write shinbun.
\end{chapteropening}

\section[\texorpdfstring{\ja{でんわ} (denwa)}{でんわ (denwa)}]{\ja{でんわ} (denwa) — a telephone}

\label{lesson:JA-C134-denwa}



\begin{quote}
Three recalls before the new word.

\begin{itemize}
\raggedright
  \item \emph{Write it:} \textbf{\ja{ぬ}} from memory — \textbf{R2}, five lessons back
  \item \emph{Write it:} \textbf{\ja{そ}} from memory — \textbf{R3}, twenty lessons back
  \item \emph{Recall:} say \emph{grey} — \textbf{R4}, eighty lessons back
\end{itemize}
\end{quote}

\subsection*{You'll want to know: \ja{でんわ}}
\textbf{\ja{でんわ}} — \emph{denwa} — \textquotedblleft{}a telephone\textquotedblright{}, on a desk or in your pocket.

Three beats: \emph{de–n–wa}. The last two signs are \textbf{\ja{ん}} and \textbf{\ja{わ}}, which you already write. The first one is \textbf{\ja{て}} with the two short strokes of the voicing mark at its upper right, so \emph{te} becomes \emph{de}. The next lesson writes it.

\subsection*{Your turn}
\begin{itemize}
\raggedright
  \item \emph{Say it:} \emph{denwa}
  \item \emph{Say it:} \emph{denwa}, then count its beats aloud — three
  \item \emph{Recall:} say which sign in \textbf{\ja{でんわ}} carries the two-stroke mark, and name the sign under it
\end{itemize}

\begin{tcolorbox}[breakable,colback=teal!4,colframe=teal!35!black,title={Before you move on}]
What is \textquotedblleft{}a telephone\textquotedblright{}? (\textbf{\ja{でんわ}}, \emph{denwa}.) How many beats? (\textbf{Three}.)
\end{tcolorbox}
\section[\texorpdfstring{\ja{で} (de)}{で (de)}]{\ja{で} — \ja{て}, voiced}

\label{lesson:JA-W134-de}



\begin{quote}
Four recalls, then the sign.

\begin{itemize}
\raggedright
  \item \emph{Recall:} say \emph{a telephone} — \textbf{R1}, one lesson back
  \item \emph{Recall:} say \emph{cloth} — \textbf{R2}, five lessons back
  \item \emph{Recall:} say \emph{outside} — \textbf{R3}, twenty lessons back
  \item \emph{Recall:} say \emph{thin}, of tea or of a colour — \textbf{R4}, eighty lessons back
\end{itemize}
\end{quote}

\subsection*{Script — the change}
\hlblockfigure{\detokenize{figures/JA-W134-de-filmstrip.pdf}}{How \ja{で} is written, stroke by stroke}

\begin{quote}
\textbf{\ja{て}} \emph{te} + \textbf{\ja{゛}} $\to$ \textbf{\ja{で}} \emph{de}
\end{quote}

Write \textbf{\ja{て}} in full: one continuous stroke, across the top, back, then down and round to the right. Then add the two short strokes of the voicing mark at the upper right, the left one first. The base shape does not change.

The mark voices the opening consonant: \emph{t} becomes \emph{d}, and the vowel stays \emph{e}. It is the change you already made on \textbf{\ja{た}} for \textbf{\ja{だ}} and on \textbf{\ja{と}} for \textbf{\ja{ど}}.

With \textbf{\ja{ん}} and \textbf{\ja{わ}} after it, you can now write the word from the last lesson: \textbf{\ja{でんわ}}, \emph{denwa}, a telephone.

\subsection*{Your turn}
\begin{enumerate}
\raggedright
  \item Write \textbf{\ja{て}}, then \textbf{\ja{で}}, and say \emph{te}, \emph{de}.
  \item Write \textbf{\ja{た} \ja{だ}}, \textbf{\ja{と} \ja{ど}} and \textbf{\ja{て} \ja{で}} in pairs, and say each one.
  \item Write \textbf{\ja{でんわ}} from memory and say \emph{denwa}.
\end{enumerate}

\begin{tcolorbox}[breakable,colback=teal!4,colframe=teal!35!black,title={Before you move on}]
Stop after one accurate form.

Source: \href{https://www.unicode.org/charts/PDF/U3040.pdf}{Unicode Hiragana chart}.
\end{tcolorbox}
\section[\texorpdfstring{\ja{でんしゃ} (densha)}{でんしゃ (densha)}]{\ja{でんしゃ} (densha) — a train}

\label{lesson:JA-C134-densha}



\begin{quote}
Four recalls before the new word.

\begin{itemize}
\raggedright
  \item \emph{Write it:} \textbf{\ja{で}} from memory — \textbf{R1}, one lesson back
  \item \emph{Recall:} say \emph{a room} — \textbf{R2}, five lessons back
  \item \emph{Recall:} say \emph{this}, a thing near you — \textbf{R3}, twenty lessons back
  \item \emph{Recall:} say \emph{shallow} — \textbf{R4}, eighty lessons back
\end{itemize}
\end{quote}

\subsection*{You'll want to know: \ja{でんしゃ}}
\textbf{\ja{でんしゃ}} — \emph{densha} — \textquotedblleft{}a train\textquotedblright{}, the kind that stops at a station.

The new sign is spent at once. \textbf{\ja{で}} and \textbf{\ja{ん}}, then \textbf{\ja{し}} with a small \textbf{\ja{ゃ}}: three beats, \emph{de–n–sha}. The small \textbf{\ja{ゃ}} joins \textbf{\ja{し}} into one beat, as it joins \textbf{\ja{ち}} in \textbf{\ja{おちゃ}}. You write every sign in it now.

\subsection*{Your turn}
\begin{itemize}
\raggedright
  \item \emph{Say it:} \emph{densha}
  \item \emph{Say it:} \emph{denwa}, then \emph{densha}, and listen for the \emph{den} in both
  \item \emph{Write it:} \textbf{\ja{でんしゃ}} from memory, with the \textbf{\ja{ゃ}} small and low
\end{itemize}

\begin{tcolorbox}[breakable,colback=teal!4,colframe=teal!35!black,title={Before you move on}]
What does \ja{でんしゃ} mean? (\textquotedblleft{}A train\textquotedblright{}.) Say it once more.
\end{tcolorbox}
\section[\texorpdfstring{\ja{かばん} (kaban)}{かばん (kaban)}]{\ja{かばん} (kaban) — a bag, a briefcase}

\label{lesson:JA-C134-kaban}



\begin{quote}
Four recalls before the new word.

\begin{itemize}
\raggedright
  \item \emph{Recall:} say \emph{a train} — \textbf{R1}, one lesson back
  \item \emph{Write it:} \textbf{\ja{へ}} from memory — \textbf{R2}, five lessons back
  \item \emph{Write it:} \textbf{\ja{れ}} from memory — \textbf{R3}, twenty lessons back
  \item \emph{Recall:} say \emph{deep} — \textbf{R4}, eighty lessons back
\end{itemize}
\end{quote}

\subsection*{You'll want to know: \ja{かばん}}
\textbf{\ja{かばん}} — \emph{kaban} — \textquotedblleft{}a bag\textquotedblright{} you carry: a school bag, a briefcase, a handbag.

Three beats: \emph{ka–ba–n}. \textbf{\ja{か}} and \textbf{\ja{ん}} you already write. The middle sign is \textbf{\ja{は}} with the voicing mark, and on an \emph{h} sign the mark makes a \emph{b}: \emph{ha} becomes \emph{ba}. The next lesson writes it.

\subsection*{Your turn}
\begin{itemize}
\raggedright
  \item \emph{Say it:} \emph{kaban}
  \item \emph{Say it:} \emph{kaban}, then count its beats aloud — three
  \item \emph{Recall:} say which sign in \textbf{\ja{かばん}} carries the two-stroke mark, and name the sign under it
\end{itemize}

\begin{tcolorbox}[breakable,colback=teal!4,colframe=teal!35!black,title={Before you move on}]
What is \textquotedblleft{}a bag\textquotedblright{} you carry to school or to work? (\textbf{\ja{かばん}}, \emph{kaban}.) How many beats? (\textbf{Three}.)
\end{tcolorbox}
\section[\texorpdfstring{\ja{ば} (ba)}{ば (ba)}]{\ja{ば} — \ja{は}, voiced to ba}

\label{lesson:JA-W134-ba}



\begin{quote}
Three recalls, then the sign.

\begin{itemize}
\raggedright
  \item \emph{Recall:} say \emph{a bag}, the one you carry to work — \textbf{R1}, one lesson back
  \item \emph{Recall:} say \emph{that}, a thing near the listener — \textbf{R3}, twenty lessons back
  \item \emph{Recall:} say \emph{precious} — \textbf{R4}, eighty lessons back
\end{itemize}
\end{quote}

\subsection*{Script — the change}
\hlblockfigure{\detokenize{figures/JA-W134-ba-filmstrip.pdf}}{How \ja{ば} is written, stroke by stroke}

\begin{quote}
\textbf{\ja{は}} \emph{ha} + \textbf{\ja{゛}} $\to$ \textbf{\ja{ば}} \emph{ba}
\end{quote}

Write \textbf{\ja{は}} in full, in the order you learned it. Then add the two short strokes of the voicing mark at the upper right, the left one first.

Here the mark does something you have met only once: on an \emph{h} sign it makes a \emph{b}, so \emph{ha} becomes \emph{ba}. You saw it on \textbf{\ja{ほ}} for \textbf{\ja{ぼ}}, in \textbf{\ja{さんぼん}}.

With \textbf{\ja{か}} in front and \textbf{\ja{ん}} after it, you can now write the word from the last lesson: \textbf{\ja{かばん}}, \emph{kaban}, a bag.

\subsection*{Your turn}
\begin{enumerate}
\raggedright
  \item Write \textbf{\ja{は}}, then \textbf{\ja{ば}}, and say \emph{ha}, \emph{ba}.
  \item Write \textbf{\ja{ほ} \ja{ぼ}} and \textbf{\ja{は} \ja{ば}} in pairs, and say each one.
  \item Write \textbf{\ja{かばん}} from memory and say \emph{kaban}.
\end{enumerate}

\begin{tcolorbox}[breakable,colback=teal!4,colframe=teal!35!black,title={Before you move on}]
Stop after one accurate form.

Source: \href{https://www.unicode.org/charts/PDF/U3040.pdf}{Unicode Hiragana chart}.
\end{tcolorbox}
\section[\texorpdfstring{\ja{ばしょ} (basho)}{ばしょ (basho)}]{\ja{ばしょ} (basho) — a place, a spot}

\label{lesson:JA-C134-basho}



\begin{quote}
Four recalls before the new word.

\begin{itemize}
\raggedright
  \item \emph{Write it:} \textbf{\ja{ば}} from memory — \textbf{R1}, one lesson back
  \item \emph{Recall:} say \emph{a telephone} — \textbf{R2}, five lessons back
  \item \emph{Recall:} say \emph{a car} — \textbf{R3}, twenty lessons back
  \item \emph{Recall:} say \emph{necessary} — \textbf{R4}, eighty lessons back
\end{itemize}
\end{quote}

\subsection*{You'll want to know: \ja{ばしょ}}
\textbf{\ja{ばしょ}} — \emph{basho} — \textquotedblleft{}a place, a spot\textquotedblright{}: where something is, or where something happens.

The new sign is spent at once. \textbf{\ja{ば}}, then \textbf{\ja{し}} with a small \textbf{\ja{ょ}}: two beats, \emph{ba–sho}. The small \textbf{\ja{ょ}} joins \textbf{\ja{し}} into one beat, the way the small \textbf{\ja{ゃ}} did in \textbf{\ja{でんしゃ}}. You write every sign in it now.

\subsection*{Your turn}
\begin{itemize}
\raggedright
  \item \emph{Say it:} \emph{basho}
  \item \emph{Say it:} \emph{kaban}, then \emph{basho}, and listen for the \emph{ba} in both
  \item \emph{Write it:} \textbf{\ja{ばしょ}} from memory, three signs for two beats
\end{itemize}

\begin{tcolorbox}[breakable,colback=teal!4,colframe=teal!35!black,title={Before you move on}]
What does \ja{ばしょ} mean? (\textquotedblleft{}A place, a spot\textquotedblright{}.) Say it once more.
\end{tcolorbox}
\section[\texorpdfstring{\ja{たべる} (taberu)}{たべる (taberu)}]{\ja{たべる} (taberu) — to eat}

\label{lesson:JA-C134-taberu}



\begin{quote}
Four recalls before the new word.

\begin{itemize}
\raggedright
  \item \emph{Recall:} say \emph{a place} — \textbf{R1}, one lesson back
  \item \emph{Write it:} \textbf{\ja{で}} from memory — \textbf{R2}, five lessons back
  \item \emph{Write it:} \textbf{\ja{る}} from memory — \textbf{R3}, twenty lessons back
  \item \emph{Recall:} say \emph{easy to do}, the opposite of complicated — \textbf{R4}, eighty lessons back
\end{itemize}
\end{quote}

\subsection*{You'll want to know: \ja{たべる}}
\textbf{\ja{たべる}} — \emph{taberu} — \textquotedblleft{}to eat\textquotedblright{}. This is the dictionary form, the one a word list gives. Said politely, it is \emph{tabemasu}.

Three beats: \emph{ta–be–ru}. \textbf{\ja{た}} and \textbf{\ja{る}} you already write. The middle sign is \textbf{\ja{へ}} with the voicing mark, so \emph{he} becomes \emph{be}. The next lesson writes it.

\subsection*{Your turn}
\begin{itemize}
\raggedright
  \item \emph{Say it:} \emph{taberu}
  \item \emph{Say it:} \emph{taberu}, then count its beats aloud — three
  \item \emph{Recall:} say which sign in \textbf{\ja{たべる}} carries the two-stroke mark, and name the sign under it
\end{itemize}

\begin{tcolorbox}[breakable,colback=teal!4,colframe=teal!35!black,title={Before you move on}]
What is \textquotedblleft{}to eat\textquotedblright{}? (\textbf{\ja{たべる}}, \emph{taberu}.) How many beats? (\textbf{Three}.)
\end{tcolorbox}
\section[\texorpdfstring{\ja{べ} (be)}{べ (be)}]{\ja{べ} — \ja{へ}, voiced to be}

\label{lesson:JA-W134-be}



\begin{quote}
Three recalls, then the sign.

\begin{itemize}
\raggedright
  \item \emph{Recall:} say \emph{to eat} — \textbf{R1}, one lesson back
  \item \emph{Recall:} say \emph{a train} — \textbf{R2}, five lessons back
  \item \emph{Recall:} say \emph{complicated} — \textbf{R4}, eighty lessons back
\end{itemize}
\end{quote}

\subsection*{Script — the change}
\hlblockfigure{\detokenize{figures/JA-W134-be-filmstrip.pdf}}{How \ja{べ} is written, stroke by stroke}

\begin{quote}
\textbf{\ja{へ}} \emph{he} + \textbf{\ja{゛}} $\to$ \textbf{\ja{べ}} \emph{be}
\end{quote}

Write \textbf{\ja{へ}} in full: one stroke, up to the low peak and down to the right. Then add the two short strokes of the voicing mark at the upper right, above the long way down, the left one first.

As with \textbf{\ja{は}} and \textbf{\ja{ば}}, the \emph{h} sign voices to \emph{b}: \emph{he} becomes \emph{be}.

With \textbf{\ja{た}} in front and \textbf{\ja{る}} after it, you can now write the word from the last lesson: \textbf{\ja{たべる}}, \emph{taberu}, to eat.

\subsection*{Your turn}
\begin{enumerate}
\raggedright
  \item Write \textbf{\ja{へ}}, then \textbf{\ja{べ}}, and say \emph{he}, \emph{be}.
  \item Write \textbf{\ja{は} \ja{ば}} and \textbf{\ja{へ} \ja{べ}} in pairs, and say each one.
  \item Write \textbf{\ja{たべる}} from memory and say \emph{taberu}.
\end{enumerate}

\begin{tcolorbox}[breakable,colback=teal!4,colframe=teal!35!black,title={Before you move on}]
Stop after one accurate form.

Source: \href{https://www.unicode.org/charts/PDF/U3040.pdf}{Unicode Hiragana chart}.
\end{tcolorbox}
\section[\texorpdfstring{\ja{べんきょう} (benkyō)}{べんきょう (benkyō)}]{\ja{べんきょう} (benkyō) — study, studying}

\label{lesson:JA-C134-benkyou}



\begin{quote}
Four recalls before the new word.

\begin{itemize}
\raggedright
  \item \emph{Write it:} \textbf{\ja{べ}} from memory — \textbf{R1}, one lesson back
  \item \emph{Recall:} say \emph{a bag}, the one you carry to work — \textbf{R2}, five lessons back
  \item \emph{Recall:} say \emph{a station} — \textbf{R3}, twenty lessons back
  \item \emph{Recall:} say \emph{well-behaved}, of a child or a pet — \textbf{R4}, eighty lessons back
\end{itemize}
\end{quote}

\subsection*{You'll want to know: \ja{べんきょう}}
\textbf{\ja{べんきょう}} — \emph{benkyō} — \textquotedblleft{}study, studying\textquotedblright{}: the work of learning, as you are doing now.

The new sign is spent at once. \textbf{\ja{べ}} and \textbf{\ja{ん}}, then \textbf{\ja{き}} with a small \textbf{\ja{ょ}}, then \textbf{\ja{う}}: four beats, \emph{be–n–kyo–o}. The \textbf{\ja{う}} at the end stretches the \emph{o}, as it does in \textbf{\ja{きのう}}. You write every sign in it now.

\subsection*{Your turn}
\begin{itemize}
\raggedright
  \item \emph{Say it:} \emph{benkyō}
  \item \emph{Say it:} \emph{taberu}, then \emph{benkyō}, and listen for the \emph{be} in both
  \item \emph{Write it:} \textbf{\ja{べんきょう}} from memory, five signs for four beats
\end{itemize}

\begin{tcolorbox}[breakable,colback=teal!4,colframe=teal!35!black,title={Before you move on}]
What does \ja{べんきょう} mean? (\textquotedblleft{}Study, studying\textquotedblright{}.) Say it once more.
\end{tcolorbox}
\section[\texorpdfstring{\ja{しんぶん} (shinbun)}{しんぶん (shinbun)}]{\ja{しんぶん} (shinbun) — a newspaper}

\label{lesson:JA-C134-shinbun}



\begin{quote}
Four recalls before the new word.

\begin{itemize}
\raggedright
  \item \emph{Recall:} say \emph{study} — \textbf{R1}, one lesson back
  \item \emph{Write it:} \textbf{\ja{ば}} from memory — \textbf{R2}, five lessons back
  \item \emph{Write it:} \textbf{\ja{き}} from memory — \textbf{R3}, twenty lessons back
  \item \emph{Recall:} say \emph{dear and missed}, of something from the past — \textbf{R4}, eighty lessons back
\end{itemize}
\end{quote}

\subsection*{You'll want to know: \ja{しんぶん}}
\textbf{\ja{しんぶん}} — \emph{shinbun} — \textquotedblleft{}a newspaper\textquotedblright{}.

Four beats: \emph{shi–n–bu–n}. \textbf{\ja{し}} and both \textbf{\ja{ん}} you already write. The third sign is \textbf{\ja{ふ}} with the voicing mark, so \emph{fu} becomes \emph{bu}. The next lesson writes it.

\subsection*{Your turn}
\begin{itemize}
\raggedright
  \item \emph{Say it:} \emph{shinbun}
  \item \emph{Say it:} \emph{shinbun}, then count its beats aloud — four
  \item \emph{Recall:} say which sign in \textbf{\ja{しんぶん}} carries the two-stroke mark, and name the sign under it
\end{itemize}

\begin{tcolorbox}[breakable,colback=teal!4,colframe=teal!35!black,title={Before you move on}]
What is \textquotedblleft{}a newspaper\textquotedblright{}? (\textbf{\ja{しんぶん}}, \emph{shinbun}.) How many beats? (\textbf{Four}.)
\end{tcolorbox}
\section[\texorpdfstring{\ja{ぶ} (bu)}{ぶ (bu)}]{\ja{ぶ} — \ja{ふ}, voiced to bu}

\label{lesson:JA-W134-bu}



\begin{quote}
Four recalls, then the sign.

\begin{itemize}
\raggedright
  \item \emph{Recall:} say \emph{a newspaper} — \textbf{R1}, one lesson back
  \item \emph{Recall:} say \emph{a place} — \textbf{R2}, five lessons back
  \item \emph{Recall:} say \emph{yesterday} — \textbf{R3}, twenty lessons back
  \item \emph{Recall:} say \emph{frustrating} — \textbf{R4}, eighty lessons back
\end{itemize}
\end{quote}

\subsection*{Script — the change}
\hlblockfigure{\detokenize{figures/JA-W134-bu-filmstrip.pdf}}{How \ja{ぶ} is written, stroke by stroke}

\begin{quote}
\textbf{\ja{ふ}} \emph{fu} + \textbf{\ja{゛}} $\to$ \textbf{\ja{ぶ}} \emph{bu}
\end{quote}

Write \textbf{\ja{ふ}} in full, in the order you learned it. Then add the two short strokes of the voicing mark at the upper right, the left one first.

The sign is read \emph{bu}: the whole \emph{h} row voices to \emph{b}, as \textbf{\ja{は}} did to \textbf{\ja{ば}} and \textbf{\ja{へ}} to \textbf{\ja{べ}}.

With \textbf{\ja{しん}} in front and \textbf{\ja{ん}} after it, you can now write the word from the last lesson: \textbf{\ja{しんぶん}}, \emph{shinbun}, a newspaper.

\subsection*{Your turn}
\begin{enumerate}
\raggedright
  \item Write \textbf{\ja{ふ}}, then \textbf{\ja{ぶ}}, and say \emph{fu}, \emph{bu}.
  \item Write \textbf{\ja{ば}}, \textbf{\ja{ぶ}}, \textbf{\ja{べ}} and \textbf{\ja{ぼ}}, and say each one.
  \item Write \textbf{\ja{しんぶん}} from memory and say \emph{shinbun}.
\end{enumerate}

\begin{tcolorbox}[breakable,colback=teal!4,colframe=teal!35!black,title={Before you move on}]
Stop after one accurate form.

Source: \href{https://www.unicode.org/charts/PDF/U3040.pdf}{Unicode Hiragana chart}.
\end{tcolorbox}
\section[(dialogue)]{Phone and train — four signs, voiced}

\label{lesson:JA-R134-phone-and-train}



\begin{quote}
Four recalls, then the review.

\begin{itemize}
\raggedright
  \item \emph{Write it:} \textbf{\ja{ぶ}} from memory — \textbf{R1}, one lesson back
  \item \emph{Recall:} say \emph{to eat} — \textbf{R2}, five lessons back
  \item \emph{Recall:} say \emph{a pond} — \textbf{R3}, twenty lessons back
  \item \emph{Recall:} say \emph{suspicious} — \textbf{R4}, eighty lessons back
\end{itemize}
\end{quote}

\subsection*{Your turn — read the seven words}
For each one below, say what it means.

Read these aloud:

\begin{itemize}
\raggedright
  \item \textbf{\ja{でんわ}} — a telephone
  \item \textbf{\ja{でんしゃ}} — a train
  \item \textbf{\ja{かばん}} — a bag
  \item \textbf{\ja{ばしょ}} — a place
  \item \textbf{\ja{たべる}} — to eat
  \item \textbf{\ja{べんきょう}} — study
  \item \textbf{\ja{しんぶん}} — a newspaper
\end{itemize}

Each of these words holds one of the four signs this chapter wrote, and each of those signs is one you already wrote, with the voicing mark added.

\subsection*{Your turn — write}
\begin{enumerate}
\raggedright
  \item \emph{Write it:} \textbf{\ja{て} \ja{で}}, \textbf{\ja{は} \ja{ば}}, \textbf{\ja{へ} \ja{べ}} and \textbf{\ja{ふ} \ja{ぶ}}, in pairs
  \item \emph{Write it:} \textbf{\ja{でんしゃ}} and \textbf{\ja{ばしょ}} from memory, and say which one you can ride
  \item Say \emph{shinbun}. \emph{Write it:} \textbf{\ja{しんぶん}}
\end{enumerate}

\begin{tcolorbox}[breakable,colback=teal!4,colframe=teal!35!black,title={Before you move on}]
A telephone? (\textbf{\ja{でんわ}}, \emph{denwa}.) A train? (\textbf{\ja{でんしゃ}}, \emph{densha}.) A bag? (\textbf{\ja{かばん}}, \emph{kaban}.) A place? (\textbf{\ja{ばしょ}}, \emph{basho}.) To eat? (\textbf{\ja{たべる}}, \emph{taberu}.) Study? (\textbf{\ja{べんきょう}}, \emph{benkyō}.) A newspaper? (\textbf{\ja{しんぶん}}, \emph{shinbun}.)
\end{tcolorbox}
